% !TEX root = ../main.tex
\section{The loss in isolation}
\label{sec:05_offpolicy_exp}

\S\ref{sec:04_bregman} is an argument about conditioning, and conditioning
arguments are cheap to make and easy to get wrong. This section measures the
five losses against each other directly.

\paragraph{Design.}
Every arm trains the same network on the same problem instances from the same
off-policy bridge anchors with the same budget; only \texttt{loss\_type}
differs. Training is off-policy throughout, deliberately: an on-policy sampler
would let the loss change the data it is subsequently trained on, and we want
the divergence isolated from the sampler. The arms are the Spence loss,
exponential-space squared error, Itakura--Saito, generalised KL
($\beta=1$), and log-space squared error.

\emph{No arm uses gradient clipping.} This is deliberate and it is what makes
the comparison a test of the tails rather than of a clipping heuristic:
Itakura--Saito and generalised KL each explode on one side
(\S\ref{sec:04_bregman}), and clipping is the obvious remedy for both. Whether
a clipped competitor closes the gap is the natural follow-up and we have not
run it, so nothing here should be read as a claim against clipped variants.

Two of these choices deserve comment. The squared-error arm is our
\emph{stabilised} implementation --- exact gradient, saturated to the float
range --- and not the naive one. Comparing against the naive version would
measure arithmetic rather than the divergence, and would flatter us: the naive
loss simply fails. The log-space arm is included as a bias probe rather than as
a candidate. It is perfectly conditioned and converges happily; by
\eqref{eq:jensen-gap} it converges to the wrong function, and the point of
including it is to show that deficit rather than to propose it.

\paragraph{Tuning, and why it matters here.}
The shared hyperparameters of the benchmark were originally fitted with the
Spence loss in place. Evaluating the alternatives at those values would be a
rigged comparison, and it is the first thing a reader should suspect. We
therefore re-tune the learning rate per (loss, dimension) over a grid spanning
$10^{-5}$ to $3\times10^{-3}$ on separate tuning seeds at a shorter budget, and
run the reported grid on fresh seeds at the chosen value. Each arm is thus
reported at its own best learning rate rather than at ours.

One caveat about that grid. At the larger dimensions several arms --- including
ours --- select the smallest rate offered, $10^{-5}$, so their optima may lie
outside the range and the tuning is truncated. It is truncated identically for
every arm, so the comparison remains matched, and it is a comparison of best
rates \emph{within} $[10^{-5},3\times10^{-3}]$ rather than of global optima. We
mention it because a boundary optimum is exactly the kind of detail that makes
a tuned-versus-tuned claim less clean than it looks.

\paragraph{Metrics.}
Two axes, because they can disagree, and the reason they can matters. The
control only ever sees $\nabla \Vf_\theta$, so a value function that is wrong
by a slowly-varying offset still steers correctly; conversely a loss can score
well on control while learning a badly biased value. The pathology of
\S\ref{sec:03_expmse} concerns the \emph{magnitude} of the learned value in
under-predicted regions, so the value axis is where it should show, and a
control-only comparison could miss it entirely.

\emph{Control quality} is the fraction of the calibrated $6$-nat headroom
closed, the same score used throughout \S\ref{sec:07_experiments}. \emph{Value quality} is the RMSE of $\Vf_\theta$
against the \emph{analytic} value over the base marginal --- available because
this problem family has $\Vf$ in closed form (App.~\ref{sec:C_experimental_details}) ---
together with its signed mean, which exposes bias rather than variance. We also
count non-finite batches, the stability axis of \S\ref{sec:03_expmse}.

% !TEX root = ../main.tex
% Auto-generated placeholder; replaced by make_ablation_table.py once the
% loss-ablation grid completes.
\begin{center}\itshape
[Results table pending: the loss-ablation grid is still running.]
\end{center}


\paragraph{What the comparison shows.}
Four readings, all paired on the same instances. The first is the one we set
out to test; the third is not the one we expected.

\emph{The Spence loss reliably repairs exponential-space squared error.} On
value RMSE it is better at every dimension we ran ($-0.37$ to $-0.75$ nats,
significant at $d=2$ and $d=32$), and on control it is ahead at every
dimension. This is the claim of \S\ref{subsec:spence} and it holds: a weight
with both tails controlled learns a more faithful value than one whose
under-prediction gradient vanishes.

\emph{It also beats generalised KL on value, everywhere.} The $\beta=1$ member
is tied with ours on control at every dimension, but its value RMSE is worse at
all four ($-0.10$ to $-0.23$ in our favour). Fixing one tail is not the same as
fixing both, and the difference shows up where the theory says it should.

\emph{But it does not dominate the family: Itakura--Saito is the better value
estimator above $d=2$.} We win that comparison at $d=2$ ($-0.12$,
significant) and lose it at $d=8$, $32$ and $128$ ($+0.11$, $+0.20$, $+0.16$;
significant at $d=128$), with control statistically tied throughout. The
direction is consistent across three dimensions, so we do not read it as
noise. A reader is entitled to conclude that on this benchmark, above two
dimensions, a divergence from 1968 estimates the value at least as well as
ours. What remains true and is worth separating out: Itakura--Saito achieves
this with a gradient that explodes for under-predictions, so it depends on the
clipping or saturation that our weight makes unnecessary; we did not clip any
arm (App.~\ref{subsec:a-numerics} saturates only where the true gradient
leaves the float range). Whether the tail control buys anything once a
competitor is properly clipped is the experiment we did not run, and it is the
one we would run next.

\emph{Conditioning and unbiasedness are different goods.} The log-space arm is
the strongest controller in the table at every dimension --- by $12$ to $26$
points --- while carrying a value bias of $1.6$ to $4.1$ nats in the direction
\eqref{eq:jensen-gap} predicts. A loss can be well conditioned and wrong. If
the value is wanted only as a gradient that trade is available and cheap; if
$\Vf$ itself is the object --- and $\Vf(0,0)=\log Z$ --- it is not a trade at
all, which is the argument of \S\ref{sec:03_expmse} made empirical.

\paragraph{One axis separates nothing.}
We counted non-finite batches per arm and found none, in any arm, at any
dimension. Once the stabilised implementations of \S\ref{sec:03_expmse} and
App.~\ref{subsec:a-numerics} are used, the overflow failures that motivated
this work do not recur, and the differences above are differences in what the
losses \emph{learn}, not in whether they run. The stability axis motivates the
numerics and is then removed by them.

\paragraph{Scope.}
This experiment isolates the divergence by construction, and therefore says
nothing about how the losses interact with an on-policy sampler, where the
loss influences its own future training data. It is also a single problem
family at one difficulty setting. The dimensional trend within it is
informative; extrapolating it to other families is not warranted.
